\documentclass{xduugthesis}
\usepackage{multirow}
\usepackage{amsmath}
\usepackage{float}
\usepackage{graphicx}
\usepackage{tabularx}
\usepackage{array}
\usepackage{xurl}

\makeatletter
\let\cleardoublepage\clearpage
\makeatother
\setlength{\emergencystretch}{2em}
\setlength{\tabcolsep}{5pt}
\renewcommand{\arraystretch}{1.12}
\newcolumntype{Y}{>{\raggedright\arraybackslash}X}

\xdusetup{
  style / font-type = font,
  style / latin-font = tacn,
  style / cjk-font = fandol,
  info / title = {基于 AI Agent 与 MCP 的手机银行语义欺诈风控系统的设计与实现},
  info / author = {云天明},
  info / department = {计算机科学与技术学院},
  info / major = {软件工程},
  info / supervisor = {谢小璞、崔笛(西安电子科技大学)},
  info / class-id = {2203054},
  info / student-id = {22009200808},
  info / abstract = {abstract-zh.tex},
  info / abstract* = {abstract-en.tex},
  info / keywords = {AI Agent；MCP协议；手机银行；语义欺诈；风控系统},
  info / keywords* = {AI Agent; MCP Protocol; Mobile Banking; Semantic Fraud; Risk Control System},
  info / bib-resource = {references.bib},
  info / acknowledgements = {acknowledgements.tex}
}

\begin{document}

\chapter{绪论}

\section{研究背景与意义}

近年来，手机银行已经从辅助渠道逐渐变成用户办理转账、账户查询、理财和缴费等业务的主要入口。移动端交易的便利性提高了金融服务效率，但也改变了风险发生的位置：很多风险不再表现为设备被攻破、账号被盗用或异地登录，而是发生在用户本人主动操作的过程中。尤其在转账场景中，诈骗者往往通过冒充客服、虚假投资、公检法威胁等话术影响用户判断，使用户在设备、账号和网络环境都看似正常的情况下完成资金转出。这类风险具有较强的语义特征和行为隐蔽性，仅依靠传统规则很难被及时识别\cite{ahi2025risks}。

传统手机银行风控系统更擅长处理结构化异常，例如交易金额突增、登录地点变化、设备指纹异常或短时间高频交易等。这类规则对于盗刷和异常登录仍然有效，但面对语义欺诈时会遇到明显局限：交易动作由用户本人发起，交易参数也可能落在正常范围内，系统很难仅凭金额和设备信息判断用户是否正在受到诱导。也就是说，风险并不一定出现在交易字段本身，而可能隐藏在用户的交易意图、输入内容和操作节奏之中。

基于这一观察，本文将研究重点放在移动端转账链路中的隐性语义欺诈识别上。一方面，借助人工智能智能体（Artificial Intelligence Agent，AI Agent）对交易备注、用户补充说明和诈骗案例进行语义分析，尝试识别传统关键词规则不易覆盖的诱导性表达；另一方面，引入输入时长、停顿次数、修改频次和页面停留等微行为特征，补充描述用户在交易过程中的操作状态。由于金融交易对安全性和可追溯性要求较高，本文并不让 AI Agent 直接执行放行或拦截动作，而是通过规则治理层将语义分析结果收敛为确定性指令。这样的设计既利用了大模型在语义理解上的优势，也尽量保留金融风控系统对稳定性、权限边界和审计链路的要求\cite{kurshan2025agentic}。

因此，本文研究具有两方面意义。理论层面上，论文尝试把认知负荷、用户微行为和语义欺诈识别结合起来，为移动端风控提供一种不同于传统静态规则的分析视角。工程层面上，论文基于 Flutter、FastAPI、AI Agent 和模型上下文协议（Model Context Protocol，MCP）实现原型系统，验证“感知、认知、治理、执行”分层架构在手机银行转账场景中的可行性，为后续智能风控系统的安全落地提供参考\cite{anthropic2024mcp}。

\section{国内外研究现状}

从国外研究和产业实践来看，金融机构较早开始使用机器学习方法处理交易风险。随机森林、XGBoost 等模型常被用于识别大额异常交易、账户盗刷和支付欺诈，其优势在于能够从历史交易数据中学习结构化风险模式\cite{bolton2002statistical,ngai2011application,dubey2025transformer}。随着生成式人工智能的发展\cite{vaswani2017attention,brown2020language}，一些研究开始关注大语言模型在金融文本处理、舆情分析、合规问答和交易描述理解中的应用。此类方法扩展了传统风控对非结构化文本的处理能力，但多数工作仍集中于通用金融文本分析，针对手机银行转账过程中“用户本人被诱导操作”的细粒度场景研究相对有限。

国内金融风控经历了从专家规则、评分卡模型到大数据用户画像的持续演进。近年来，围绕大模型轻量化部署、欺诈话术识别、多模态风险检测等方向已经出现较多探索\cite{nie2025multimodal}。这些研究能够在一定程度上识别诈骗文本或异常交易模式，但如果直接应用到手机银行核心交易场景，仍会遇到几个问题：一是模型输出具有概率性，难以直接作为资金交易的最终裁决；二是不同业务系统的数据接口不统一，影响风险上下文的完整获取；三是很多方案强调识别精度，却对审计追踪、权限隔离和人工复核等金融合规要求讨论不足\cite{anthropic2024mcp,wang2025compliance}。

AI Agent 技术为复杂任务拆解、上下文管理和工具调用提供了新的实现方式\cite{agnihotram2025agentic}。在智能办公、问答系统和自动化流程中，Agent 已经能够根据任务目标调用外部工具并生成结构化结果。但在金融风控场景中，Agent 的角色需要受到更严格的限制。本文认为，Agent 更适合作为认知分析组件，而不应直接进入资金执行链路。原因在于大模型仍可能出现误判、幻觉或输出不稳定，如果缺少规则治理和审计约束，可能带来新的业务风险。

用户微行为分析是另一个与本文相关的研究方向。已有研究表明，用户在压力、诱导或高认知负荷状态下，输入节奏、停顿时间和修改行为可能发生变化。这类特征过去更多用于用户体验优化、身份识别或异常行为检测，在实时金融风控中的应用还不充分。本文将该方向与语义欺诈识别结合\cite{li2024cognitive}，尝试用移动端微行为补充交易文本和结构化字段的不足。

MCP 协议则为智能体与外部资源之间的标准化交互提供了新的接口思路\cite{anthropic2024mcp}。与传统自定义接口相比，MCP 更强调资源、工具和提示词的统一组织方式，便于在复杂系统中约束智能体的访问边界。现有应用多集中在通用智能体中台或开发者工具场景，面向银行强监管业务的定制化实践仍较少\cite{anthropic2024mcp}。基于此，本文将 MCP 引入手机银行风控原型系统，用于封装风险计算、转账执行和审计记录等工具，并配合权限隔离与 Trace ID 追踪机制，探索其在金融级审计链路中的应用价值。

\section{现有工作不足与本文区别}

结合国内外研究现状可以看出，现有金融风控研究已经在结构化异常检测、机器学习欺诈识别、金融文本理解和大模型辅助分析等方面取得了一定进展。但面向手机银行转账场景时，现有工作仍难以完整覆盖“用户本人被诱导操作”的语义欺诈风险。该类风险的特殊性在于，交易往往由用户本人在常用设备和正常网络环境下主动发起，交易金额、登录地点和设备指纹等字段未必出现明显异常，风险更多隐藏在用户转账意图、交易备注、补充说明和操作节奏之中。

首先，传统规则和结构化模型难以覆盖语义欺诈中的真实交易意图。现有风控系统通常依据交易金额、设备环境、登录地点、收款账户和交易频次等字段进行判定，这些方法对于盗刷、撞库、异地登录和异常高频交易具有较好效果。但是在社交工程类欺诈中，诈骗者并不一定攻击账户或设备，而是通过冒充客服、虚假投资、公检法威胁等话术影响用户判断。此时，系统看到的是一次由用户本人确认的正常交易，而真正的风险隐藏在“为什么转账”和“用户是否被诱导”之中。单纯依靠硬编码规则，很难捕捉这类语义层面的风险\cite{nicolaou2025llm}。

其次，已有大语言模型或 AI Agent 方案虽然增强了文本理解能力，但并不适合直接控制金融交易执行。大模型能够识别诱导性表达、归纳诈骗话术并生成风险解释，但模型输出具有概率性，可能受到提示词、上下文覆盖范围和推理随机性的影响。如果将模型输出直接作为放行或拦截依据，一旦发生误判，就可能造成用户体验损害或资金安全风险。因此，在金融风控场景中，Agent 更适合作为认知分析组件，而不应直接进入资金执行链路。本文将 Agent 的职责限定为语义分析、风险线索提取和报告说明，最终处置仍由规则治理层根据阈值、状态机和审计约束生成确定性指令。

再次，普通接口调用式的 Agent 系统缺少适合金融场景的工具边界和审计约束。许多智能体应用通过自定义 API 调用外部工具，但这种方式在银行转账链路中容易带来权限不清、调用不可追踪、重复请求难以幂等控制等问题。对于金融业务而言，系统不仅要回答“模型判断了什么”，还要回答“模型访问了哪些数据、调用了哪些工具、工具返回了什么结果、最终状态为什么发生变化”。因此，本文引入 MCP 协议，将风险计算、审计写入、交易状态变更和报告生成等能力封装为受控工具，并通过 Transfer ID、Trace ID 和幂等校验记录每一次调用。MCP 在本文中的作用不是直接提高模型识别率，而是为 AI Agent 接入风控流程提供安全数据中枢、工具调用边界和审计追踪机制\cite{anthropic2024mcp}。

最后，现有部分研究更关注单点识别效果，而对风险处置闭环讨论不足。手机银行风控并不止于输出一个风险分数，还需要支持预检、二次核验、风险报告、用户提示、人工复核和审计留痕等流程。如果缺少状态机约束和复核机制，系统即使能够识别部分风险，也难以满足金融业务对可解释、可追溯和可复核的要求。基于此，本文设计 HIRD 分层架构，将感知、认知、治理和执行拆分为相对独立的模块：感知层采集交易文本和微行为特征，认知层由 Agent 完成语义分析，规则治理层将概率性分析结果收敛为确定性指令，业务执行层通过 MCP 工具链完成受控执行和审计记录。

综上，本文与现有工作的区别不在于单独提出某一种新的欺诈识别模型，而在于面向手机银行转账这一强安全、强审计场景，将微行为感知、AI Agent 语义分析、规则治理、MCP 工具封装、Trace ID 审计追踪和人工复核机制组合成一个可运行的风控原型闭环。该设计既利用了大模型在语义理解方面的优势，又通过 MCP 和规则治理层限制模型权限，使系统能够在测试环境中验证 AI Agent 安全接入手机银行风控流程的可行性。

\section{本文研究内容与组织结构}

\subsection{主要研究内容}

围绕上述问题，本文的研究内容主要集中在手机银行转账业务中的智能风控原型设计与实现。本文首先从交易链路出发，梳理正常转账、二次核验、风险拦截和人工复核等关键流程，明确系统需要处理的数据来源和权限边界。在此基础上，构建 HIRD 四层架构，将感知、认知、治理和执行拆分为相对独立的模块\cite{agnihotram2025agentic}。

在感知层，本文基于 Flutter 客户端实现微行为采集，重点记录输入时长、停顿次数、修改频次、页面停留和操作节奏方差等特征，并将其与交易金额、收款信息和交易备注一起封装为风险快照。在认知层，本文引入 AI Agent 和 RAG 检索机制，对交易文本、用户补充说明和风险案例进行语义匹配，输出风险线索和报告摘要。在规则治理层，本文设计动态权重与分级阈值机制，将 Agent 的概率性分析结果与硬规则评分进行整合，生成放行、二次核验或拦截等确定性指令。在业务执行层，本文通过 MCP 协议封装风险计算、转账执行和审计记录等工具，并保留 Trace ID、幂等校验和人工复核入口，使系统能够在测试环境中完成从感知采集到审计留痕的闭环。

\subsection{论文组织结构}

全文共分为六章，章节安排如下。

第一章为绪论，主要说明本文选题来源、研究背景、国内外研究现状、现有工作不足与本文区别，并给出本文的研究内容和整体组织方式。

第二章为相关技术综述，围绕本文系统实现所需的关键技术展开，介绍大语言模型与 AI Agent、MCP 协议、认知负荷理论、RAG 检索增强生成、Flutter 和 FastAPI 等内容，为后续架构设计提供技术基础\cite{sommerville2016software,pressman2019software}。

第三章为系统需求分析与总体设计，结合手机银行转账场景拆解功能性和非功能性需求，进一步给出 HIRD 四层架构、微行为特征建模、Agent 协作流程、动态权重算法和审计链路设计。

第四章为系统实现与功能验证，说明前端、后端和 MCP 工具链的工程实现方式，并结合实际界面和接口流程展示系统如何完成风险预检、二次质询、报告生成和人工复核\cite{pressman2019software}。

第五章为实验测试与分析，基于自动化测试脚本和场景化样本，对系统功能、响应时延、场景识别结果、影子模式一致性和消融实验结果进行分析。

第六章为结论与展望，总结本文完成的主要工作和实验结论，说明当前原型系统仍存在的数据规模、场景覆盖和消融验证不足，并提出后续改进方向。


\chapter{相关技术综述}
围绕本文构建的 HIRD（Harvesting-Intelligence-Regulation-Delivery，感知-认知-治理-执行）架构，本章介绍系统实现中实际用到的几类技术。与单纯罗列概念不同，本章更关注这些技术在手机银行风控原型中的作用边界：大语言模型与 AI Agent 负责语义分析，MCP 协议负责工具封装和调用约束，认知负荷理论为微行为特征提供解释依据，RAG、Flutter 和 FastAPI 则分别支撑知识检索、移动端采集和后端异步编排。

\section{大语言模型与金融 AI Agent 理论}
大语言模型（Large Language Model，LLM）通常基于 Transformer 架构，通过大规模语料训练获得较强的自然语言理解和生成能力\cite{nicolaou2025llm}。在手机银行风控场景中，它的价值主要体现在两点：一是能够处理交易备注、聊天片段、用户补充说明等非结构化文本；二是能够对同一类诈骗话术的变体进行语义层面的归纳，而不完全依赖固定关键词。例如，虚假理财、冒充客服和公检法威胁等场景中，诈骗者的具体措辞可能不断变化，但其诱导用户转账的意图往往具有相似结构\cite{wang2025fraud}。

不过，大语言模型并不适合直接作为金融交易的最终裁决器。模型输出可能受到提示词、上下文长度、案例覆盖程度和推理随机性的影响，存在误判或解释不稳定的风险。因此，本文在使用大模型时只把它放在认知分析环节，用于识别风险线索、生成语义摘要和辅助报告说明，而不让其直接执行放行、拦截或扣款等业务动作。

AI Agent 可以理解为在大语言模型基础上增加任务编排、上下文管理和工具调用能力的应用结构\cite{zangana2025llm}。与单次模型问答相比，Agent 更适合处理需要多步判断的任务，例如先读取交易上下文，再匹配风险案例，最后输出结构化风险原因。在本文系统中，Agent 主要位于 HIRD 架构的认知决策层，接收感知层生成的风险快照，并调用知识库检索、语义分析和报告生成等能力。

为了适配金融场景，本文对 Agent 的权限边界作了限制：Agent 可以分析、归纳和解释，但不能直接调用资金执行工具。系统最终的业务动作由规则治理层和业务执行层完成。这种设计使 Agent 能够发挥语义理解优势，同时避免模型输出直接影响交易状态。

\section{MCP 深度解析}
模型上下文协议（Model Context Protocol，MCP）是一种面向智能体应用的上下文交互协议，主要用于规范模型、工具和外部资源之间的连接方式。对于本文系统而言，MCP 的作用不是提高模型本身的识别能力，而是为“Agent 能访问什么、能调用什么、调用后如何记录”提供统一约束。手机银行风控涉及交易数据、用户行为数据和审计日志，如果仍采用零散接口直接拼接，后续排查和权限管理都会变得困难。

MCP 协议中的 Resources、Tools 和 Prompts 三类机制分别对应不同问题。Resources 用于描述可访问的资源，例如风险案例库、交易上下文或审计记录；Tools 用于封装具体操作，例如风险评分、转账执行、个人身份信息（Personally Identifiable Information，PII）脱敏和日志落盘；Prompts 则用于约束 Agent 的输入输出格式，使语义分析结果更容易被后续规则治理层解析。通过这种方式，系统可以把复杂能力拆成边界清晰的原子工具，而不是让 Agent 直接接触底层业务逻辑。

在本文实现中，MCP 更接近一个安全数据中枢。Agent 需要通过 MCP 工具获取必要上下文，工具内部再完成权限检查、幂等校验和 Trace ID 记录。这样做的好处是，即便后续更换模型或调整提示词，底层交易执行和审计链路仍由固定工具控制。对金融风控系统来说，这种可追踪、可复核的调用方式比单纯追求模型输出更重要。

\section{认知负荷理论与微行为特征分析}
认知负荷理论认为，人在处理信息时会受到注意力、短时记忆和任务复杂度的限制\cite{sweller1988cognitive}。当用户在转账过程中同时面对诈骗者话术、时间压力、身份伪装或威胁暗示时，其认知资源会被额外占用，操作行为也可能出现变化\cite{ibraheem2025duress}。这类变化不一定会体现在交易金额或设备信息中，但可能反映在输入速度、停顿位置、修改次数和页面停留时间等细节上。因此，在语义欺诈识别中，仅分析交易字段是不够的，还需要观察用户完成交易动作时的过程特征。

本文没有把微行为特征直接等同于欺诈结论，而是将其作为辅助风险证据使用。原因在于，用户输入变慢或频繁修改并不一定意味着被诈骗，也可能来自网络卡顿、年龄差异、输入习惯或临时犹豫。因此，本文在感知层采集有效输入时长、非预期停顿次数、内容修改频次、页面停留时长和操作节奏方差等指标后，并不单独据此拦截交易，而是将其与交易文本、设备状态和规则评分共同送入后续风控链路。

从系统设计角度看，微行为特征的价值在于补充传统结构化字段的盲区。对于“用户本人主动转账”的社交工程欺诈，设备和账号可能完全正常，金额也可能被控制在规则阈值以内。此时，用户在交易界面上的停顿、反复修改和异常节奏，能够为认知决策层提供额外线索。本文后续在 HIRD 架构中将这类数据封装为感知快照，使 Agent 和规则治理层能够在同一条风险链路中使用这些信息。

需要注意的是，微行为特征存在较强的个体差异，不能简单套用统一阈值。本文原型系统主要通过场景化样本验证该思路的可行性，后续若接入真实业务环境，还需要结合用户历史行为基线、设备性能差异和更多异常样本进行校准。

\section{检索增强生成（Retrieval-Augmented Generation，RAG）与知识库构建}
检索增强生成（RAG）的基本思路，是在大语言模型生成回答前，先从外部知识库中检索与当前问题相关的材料，再把检索结果作为上下文提供给模型\cite{lewis2020rag,srinivasan2025rag}。在本文场景中，RAG 的作用不是让模型“知道更多金融知识”这么简单，而是让 Agent 的风险判断能够尽量依托具体案例和规则片段，减少仅凭模型内部知识进行推断带来的不稳定性。

语义欺诈具有明显的变体特征。同一类诈骗可能使用不同表达，例如虚假理财可以表现为“内部收益通道”“限时认购”“稳赚回款”等多种话术；冒充客服也可能通过订单异常、退款验证、账户冻结等不同理由诱导转账。单纯依靠关键词匹配时，规则库需要不断枚举具体词语，维护成本较高。RAG 结合向量检索后，可以从语义相似度角度匹配历史案例，使系统对同类变体话术具备一定的泛化能力。

本文知识库主要围绕三类内容组织：一是典型欺诈案例，用于提供场景匹配依据；二是风控规则说明，用于约束 Agent 输出的风险类型和处置建议；三是系统报告模板，用于规范风险摘要和风险因子字段。Agent 在处理交易备注或用户补充说明时，会先检索相似案例，再结合当前交易金额、收款信息和微行为快照进行判断。这样生成的风险说明更容易追溯到具体案例，也方便后续人工复核。

同时，RAG 也存在局限。检索结果的质量依赖知识库覆盖范围和向量匹配效果，如果案例库过小或样本过于模板化，Agent 仍可能遗漏新型诈骗话术。本文当前知识库主要服务于原型系统和场景化测试，后续需要通过更丰富的真实或半真实案例持续补充，才能支撑更复杂的生产环境。

\section{跨平台开发与异步编排技术}
本文原型系统选择 Flutter 作为移动端实现框架，主要考虑其跨平台能力和事件监听机制。手机银行风控采集不能明显打断用户操作，也不能让用户为了测试而额外输入数据。因此，感知层需要在转账流程中记录输入时长、停顿次数、修改频次和页面停留等行为数据，并在用户提交交易时生成统一快照。Flutter 的组件化界面和事件回调机制适合承载这种采集逻辑，也便于在同一套代码中模拟不同移动端交互流程。

在本文实现中，移动端采集不是为了记录用户隐私内容，而是为了形成可计算的行为指标。系统关注的是“输入用了多久”“是否出现非预期停顿”“是否反复修改”等过程信息，而不是简单保存用户完整输入过程。这样做可以在保留风险判断线索的同时，减少不必要的敏感数据暴露。后续与 MCP 数据中枢结合时，这些指标会和交易文本、交易元数据一起进入风险快照。

后端部分采用 FastAPI，主要原因是风控链路中存在多类耗时差异明显的任务：预检热路径需要在 100ms 左右完成，二次核验和 Agent 语义分析则可能受到大模型接口、网络请求和检索过程影响。FastAPI 基于异步机制组织接口调用，便于把本地规则判断、日志记录、模型调用和状态机流转拆开处理。在第五章测试中，预检路径和二次核验路径的时延差异也说明，系统需要通过双路径方式处理不同复杂度的交易请求。

因此，Flutter 和 FastAPI 在本文中并不是单纯的前后端技术选型，而是分别服务于“无感采集”和“异步风控编排”两个目标。前者负责在移动端捕捉交易过程，后者负责在后端把规则、Agent、MCP 工具和审计日志组织为可追踪流程。

\section{本章小结}
本章围绕 HIRD 架构所需的关键技术进行了说明。大语言模型和 AI Agent 为交易文本、用户说明和欺诈案例之间的语义分析提供能力，但在本文系统中只承担认知分析职责，不直接控制业务执行。MCP 协议用于封装资源和工具调用，使 Agent 的访问边界、执行动作和审计记录能够被统一约束。认知负荷理论和微行为特征分析为移动端操作过程提供解释依据，使系统能够从用户输入节奏和停顿行为中获得补充风险线索。

在工程实现方面，RAG 与知识库用于增强 Agent 对欺诈案例的检索和引用能力，降低单纯依赖模型内部知识带来的不稳定性；Flutter 支撑移动端无感采集，FastAPI 支撑后端异步编排和双路径决策。上述技术共同服务于本文的核心目标：在不让 AI 直接接管资金执行的前提下，将语义理解、行为感知和规则治理纳入同一条可审计的手机银行风控链路。后续第三章将在这些技术基础上展开系统需求分析和 HIRD 架构设计。

\chapter{系统需求分析与总体设计}

前两章已经说明，本文关注的不是一般意义上的交易异常检测，而是手机银行转账过程中较难被传统规则识别的语义欺诈。围绕这一场景，本章从实际业务链路出发，先拆解系统需要满足的功能与约束，再给出 HIRD 分层架构、感知特征建模和后续决策链路设计。章节中的需求并非独立罗列，而是服务于同一目标：在不让 AI 直接控制资金执行的前提下，把移动端行为、交易语义和规则治理组织成可测试、可审计的风控流程。

\section{系统业务需求分析}

本文原型系统以手机银行转账为主要业务对象。一次完整交易通常包含用户填写信息、系统预检、必要时触发二次核验、最终放行或拦截，以及后续查看风险报告和人工复核等步骤。因此，需求分析需要同时覆盖前端采集、后端判断、AI 分析、规则收敛和审计留痕，而不能只关注某一个识别模型的准确率。

\subsection{功能性需求}

结合本文系统实现，功能性需求被划分为五类：多模态风险感知、语义欺诈识别、安全数据交互、动态分级治理和全链路审计。前两类需求负责发现风险，中间两类需求负责约束 AI 分析和生成确定性处置，最后一类需求用于支撑事后复核。

\textbf{F-01 多模态风险感知功能}

在社交工程类欺诈中，用户往往使用本人设备和本人账号发起交易，金额、IP 地址等字段也可能没有明显异常。为避免系统只看到“参数正常”的一面，感知模块需要在不打断用户转账流程的前提下记录操作过程，例如停顿、修改、页面停留和设备上下文，并将这些信息与交易参数放入同一份风险快照中。

\begin{table}[H]
    \centering
    \caption{多模态风险感知功能细分}
    \label{tab:f01-sub}
    \begin{tabularx}{\linewidth}{|c|Y|}
        \hline
        细分功能名称 & 核心任务说明 \\
        \hline
        异步非阻塞行为监听 & 利用前端框架在后台静默采集按键停顿等数据，不干扰转账界面正常交互 \\
        \hline
        微行为数据量化 & 将输入停顿时间、反复修改动作等操作特征转换为系统可计算的数值指标 \\
        \hline
        交易环境特征整合 & 同步提取登录位置、设备型号等基础结构化信息，与行为数据对齐 \\
        \hline
        风险感知快照生成 & 将多源数据按时间线对齐，封装为固定格式的快照供下游模块消费 \\
        \hline
    \end{tabularx}
\end{table}

\textbf{F-02 语义欺诈智能识别功能}

很多诈骗话术并不会稳定使用同一组关键词。即使都属于虚假理财或冒充客服场景，用户看到的表达也可能不断变化。因此，系统需要能够分析交易备注、用户补充说明和案例库中的语义关系，识别隐藏在自然语言背后的诱导意图。该功能输出的结果只作为风险线索和评分输入，不直接转化为业务执行动作。

\begin{table}[H]
    \centering
    \caption{语义欺诈智能识别功能细分}
    \label{tab:f02-sub}
    \begin{tabularx}{\linewidth}{|c|Y|}
        \hline
        细分功能名称 & 核心任务说明 \\
        \hline
        复杂话术意图提取 & 解析转账附言或聊天记录中的诱导性文本，识别其真实交易目的 \\
        \hline
        知识库相似度比对 & 将文本特征与 RAG 知识库中的历史欺诈案例进行语义级相似度匹配 \\
        \hline
        隐性风险分值计算 & 综合上下文语义理解与案例匹配结果，输出量化的语义风险评分 \\
        \hline
    \end{tabularx}
\end{table}

\textbf{F-03 安全数据交互与权限管控功能}

Agent 在分析交易风险时需要读取上下文，但银行业务数据具有较高敏感性，不能简单把原始交易信息直接暴露给模型。本功能要求系统在接口层面限制 Agent 的访问范围，只提供完成语义分析所需的脱敏特征和必要上下文，并通过 MCP 工具封装方式记录每一次调用。

\begin{table}[H]
    \centering
    \caption{安全数据交互与权限管控功能细分}
    \label{tab:f03-sub}
    \begin{tabularx}{\linewidth}{|c|Y|}
        \hline
        细分功能名称 & 核心任务说明 \\
        \hline
        数据访问权限隔离 & 限制大模型仅能读取脱敏后的特征数据，禁止直接访问底层数据库 \\
        \hline
        原子化风控工具封装 & 基于 MCP 协议将风险计算等动作封装为独立工具，规范调用路径 \\
        \hline
        敏感隐私信息脱敏 & 在数据传递给外部模型前，自动对姓名、卡号等字段进行掩码处理 \\
        \hline
    \end{tabularx}
\end{table}

\textbf{F-04 动态分级风控决策与治理功能}

大模型输出更适合描述风险可能性，而金融交易需要明确动作。因此，系统需要在治理层把语义评分、规则评分和场景匹配程度重新整合，输出放行、二次核验或拦截等确定性结果。该功能的重点不是让 AI 取代规则，而是在规则框架内吸收语义分析的补充价值。

\begin{table}[H]
    \centering
    \caption{动态分级风控决策与治理功能细分}
    \label{tab:f04-sub}
    \begin{tabularx}{\linewidth}{|c|Y|}
        \hline
        细分功能名称 & 核心任务说明 \\
        \hline
        场景匹配度计算 & 评估当前交易与已知诈骗模式的相似程度，决定两类评分的权重分配 \\
        \hline
        动态权重参数调节 & 当交易呈现异常特征时，自动提升 AI 评分的权重占比进行融合计算 \\
        \hline
        综合风险等级划分 & 依据融合后的总分值，将交易判定为高、中、低三个风险等级 \\
        \hline
        分级路由指令下发 & 根据风险等级自动执行放行、人脸核验或交易拦截等差异化处置 \\
        \hline
    \end{tabularx}
\end{table}

\textbf{F-05 全链路风控审计溯源功能}

对于被拦截或进入二次核验的交易，系统需要保留足够的过程信息，方便用户查看、后台复核和后续问题定位。本功能要求将感知快照、语义分析结果、权重计算过程、最终指令和人工复核状态纳入同一条审计链路，避免只保存一个简单的“拦截”结果。

\begin{table}[H]
    \centering
    \caption{全链路风控审计溯源功能细分}
    \label{tab:f05-sub}
    \begin{tabularx}{\linewidth}{|c|Y|}
        \hline
        细分功能名称 & 核心任务说明 \\
        \hline
        原始输入快照存档 & 完整保存系统采集的行为参数与设备信息，作为事后核查的底层依据 \\
        \hline
        决策推导过程记录 & 将大模型话术分析、案例匹配情况及权重分配逻辑写入审计日志 \\
        \hline
        处置指令存储 & 将最终指令及操作时间持久化存储，确保记录不可篡改 \\
        \hline
    \end{tabularx}
\end{table}

\subsection{非功能性需求}

非功能性需求决定了系统能否真正放入转账链路中运行。本文主要从实时性、安全性、准确性和可扩展性四个方面给出约束指标，如表\ref{tab:nonfunctional-req}所示。这些指标后续将在第五章通过自动化测试、场景样本和性能统计进行验证。

\begin{table}[H]
    \centering
    \caption{系统非功能性需求}
    \label{tab:nonfunctional-req}
    \begin{tabularx}{0.75\linewidth}{|c|c|Y|}
        \hline
        需求编号 & 指标项 & 具体要求 \\
        \hline
        NF-01 & 实时性 & 预检 $\leq 100\mathrm{ms}$，含 AI 链路 $\leq 3\mathrm{s}$ \\
        \hline
        NF-02 & 安全性 & 敏感信息全脱敏，接口幂等校验 \\
        \hline
        NF-03 & 准确性 & 欺诈召回率 $\geq 90\%$，正常交易直接误拦截率 $\leq 10\%$ \\
        \hline
        NF-04 & 可扩展性 & 层级解耦，可独立升级与扩充 \\
        \hline
    \end{tabularx}
\end{table}

\textbf{NF-01 实时性}

转账预检处在用户提交交易后的关键节点，如果等待时间过长，会直接影响手机银行的交互体验。因此，本文将预检热路径控制在 $100\mathrm{ms}$ 以内，该路径只执行本地规则、结构化字段判断和必要的状态记录。当交易进入二次核验或 Agent 语义分析时，系统允许更长的异步处理时间，并通过前端等待状态和后端状态机避免阻塞主流程。

\textbf{NF-02 安全性}

安全性主要体现在三个方面：敏感信息在进入 Agent 分析前需要脱敏；接口调用需要进行权限校验；涉及交易执行或状态变更的 MCP 工具需要具备幂等保护。这样即使出现网络重试、重复提交或并发调用，也不会造成重复扣款、重复审计或交易状态漂移。

\textbf{NF-03 准确性}

准确性需要同时考虑欺诈召回和正常交易直接误拦截。本文希望系统能够尽量识别典型语义欺诈样本，同时避免对正常交易造成过度干预。由于当前实验数据主要为场景化样本，本文将该指标作为原型系统测试目标，并在第五章中结合 54 个核心场景样本进行统计分析。

\textbf{NF-04 可扩展性}

可扩展性主要来自分层解耦。感知层可以增加新的行为指标，认知层可以替换模型或扩充知识库，治理层可以调整权重和阈值，执行层可以增加新的复核或提示方式。各层通过标准化快照、MCP 工具和明确状态字段衔接，尽量减少局部修改对整体系统的影响。

综合来看，上述约束之间存在明显取舍。实时性要求系统不能在每一笔交易中都调用大模型；识别能力又要求系统在规则难以判断时引入语义分析；安全性则要求 Agent 不能直接访问底层业务数据或执行交易动作。因此，本文采用热路径与语义路径分流的方式：常规交易优先由本地规则快速处理，疑似风险交易再进入二次核验和 Agent 分析；同时通过 MCP 工具和规则治理层把分析权限与执行权限分开。这一取舍构成了 HIRD 架构的设计出发点。

\section{系统 Use-case 建模与关键用例规格}

为进一步明确系统功能边界和参与角色，本文在业务需求分析基础上补充 Use-case 建模。本文系统主要面向手机银行转账过程中的语义欺诈风险识别，参与者包括普通用户、审核人员、风控后台服务、AI Agent 服务和 MCP 工具层。其中，普通用户负责发起转账、完成二次核验、查看风险报告和提交人工复核；审核人员负责处理复核申请；风控后台服务负责预检、状态流转和处置指令生成；AI Agent 服务负责语义风险分析；MCP 工具层负责受控工具调用、幂等校验和审计记录\cite{larman2004applying}。

\begin{figure}
    \centering
    \includegraphics[width=1.0\linewidth]{project_51_43215b44.png}
    \caption{系统总体 Use-case 图}
    \label{fig:placeholder}
\end{figure}

系统总体 Use-case 图用于展示普通用户、审核人员、风控后台服务、AI Agent 服务和 MCP 工具层之间的用例关系。该图强调的是系统参与者与业务动作之间的边界：普通用户只通过移动端发起交易、查看提示和提交复核；审核人员只处理复核单；风控后台服务负责交易状态和处置指令；AI Agent 服务只输出语义分析结果；MCP 工具层则负责工具调用边界、幂等校验和审计写入。系统并不是由 AI Agent 直接完成交易处置，而是将 Agent 的语义分析结果交由风控后台服务和规则治理流程进一步收敛。

\begin{table}[H]
    \centering
    \caption{关键 Use-case 规格说明}
    \label{tab:key-use-case-spec}
    \begin{tabularx}{\linewidth}{|c|Y|}
        \hline
        用例名称 & 主要参与者、前置条件、基本流程与后置结果 \\
        \hline
        普通用户发起转账 & 普通用户。用户已登录并进入转账页面。用户填写收款方、金额和备注后提交交易；系统生成交易流水号并进入风险预检流程。 \\
        \hline
        系统执行风险预检 & 风控后台服务、MCP 工具层。系统已收到交易请求和感知快照。后台读取交易参数、设备信息、行为特征和备注文本，计算风险结果并输出放行、二次核验或拦截建议。 \\
        \hline
        用户完成二次核验 & 普通用户、AI Agent 服务。预检结果为需要进一步确认。用户补充说明转账原因或完成核验；系统重新分析语义风险并更新交易处置状态。 \\
        \hline
        系统生成风险报告 & 风控后台服务、AI Agent 服务。交易进入风险处置流程。系统汇总风险等级、触发因子、语义线索和处置依据，形成用户和审核人员均可查看的风险报告。 \\
        \hline
        用户提交人工复核 & 普通用户。用户对系统处置结果存在异议。用户提交复核申请；系统创建复核单，并关联原交易流水、风险报告和审计记录。 \\
        \hline
        审核人员处理复核单 & 审核人员。复核单处于待处理状态。审核人员查看交易上下文、系统报告和审计链路，给出维持、解除或补充核查等处理结论。 \\
        \hline
    \end{tabularx}
\end{table}

上述 Use-case 规格说明补充了系统从用户提交转账到后台复核处理的主要业务路径。其中，“系统执行风险预检”和“用户完成二次核验”对应自动化风控流程，“系统生成风险报告”和“审核人员处理复核单”对应解释与复核流程。通过这种方式，第三章的功能需求不仅停留在模块说明层面，也进一步明确了系统参与者、触发条件和业务结果。

\section{HIRD 系统总体架构设计}

根据前述需求拆解，本文提出 HIRD 四层分层风控架构。HIRD 分别对应感知集成层、认知决策层、规则治理层和业务执行层。该架构的核心思路是让不同能力停留在各自边界内：感知层只负责形成数据，认知层只负责分析风险，治理层负责把分析结果转成明确指令，执行层只执行已经确定的业务动作。这样设计的目的，是避免 AI 分析结果绕过规则治理直接影响资金流转。

\subsection{架构层级定义与设计理念}

本文将系统拆解为四个层级，并不是为了增加架构复杂度，而是为了明确每一类能力的边界\cite{fowler2002patterns}。感知集成层（H 层）负责采集交易文本、结构化参数和微行为数据；认知决策层（I 层）负责语义分析、案例匹配和风险线索归纳；规则治理层（R 层）负责融合规则评分和语义评分，输出确定性处置；业务执行层（D 层）负责前端提示、交易放行、拦截、报告展示和人工复核入口。架构总览如图 3.1 所示。

\begin{figure}
    \centering
    \includegraphics[width=1.0\linewidth]{project_51_50362823.png}
    \caption{HIRD四层风控架构图}
    \label{fig:placeholder}
\end{figure}

上述四层划分的关键，是把“分析权限”和“执行权限”分开。I 层可以读取脱敏后的上下文并生成风险线索，但不能直接调用转账执行或拦截工具；R 层根据规则和阈值决定最终指令；D 层只接收明确指令并完成页面反馈或状态变更。这样即使 Agent 的语义判断出现偏差，也不会直接越过治理层影响资金动作。

\subsection{层间数据流转与协作机制}

一笔交易从用户提交到最终完成，会在 HIRD 四层之间按固定方向流转。系统使用全局唯一交易流水号贯穿感知快照、语义分析、治理决策和审计记录，确保每个环节都能被定位。根据预检结果不同，交易可能走快速放行或拦截路径，也可能进入二次质询、报告生成和人工复核路径。

\begin{figure}
    \centering
    \includegraphics[width=1.0\linewidth]{project_51_4adf40b6.png}
    \caption{各层间数据流转}
    \label{fig:placeholder}
\end{figure}

预检路径覆盖所有交易的首轮判定。用户提交转账请求后，H 层生成包含交易参数、设备上下文、文本内容和微行为指标的感知快照。快照进入 I 层后，系统对交易文本进行基础分类和案例匹配，形成风险类别、匹配场景和语义线索。R 层再结合结构化规则、行为特征和语义结果计算综合风险分值，并映射为 pass、block 或 interrogate 三类结果。pass 表示交易可继续执行，block 表示直接拦截，interrogate 表示需要进入二次质询。整个预检过程的关键节点写入审计日志，前端同步展示决策结果和必要解释信息。

二次质询路径用于处理预检阶段难以直接放行或拦截的交易。R 层先检查交易状态是否允许进入二次质询，随后由前端提示用户补充说明。I 层的 Agent 服务会先根据本地红旗规则判断是否存在明确恶意回复，再判断是否属于低风险合规回复；当本地规则无法覆盖时，再调用大模型进行语义分析。分析结果回到 R 层后，R 层更新交易状态、写入审计日志，并触发风险报告生成。前端随后获得二次决策结果、语义红旗标识和解释面板数据。

风险报告查看与人工复核路径用于补充自动化风控的解释和兜底能力。用户可以在前端查看风险标题、综合风险等级、风险摘要和风险因子。如果用户对拦截或核验结果有异议，可以从报告详情页发起复核申请。系统创建复核记录后进入后台处理队列，审核人员可查看交易上下文、审计链路和系统解释，再给出处置结论。复核结果作为治理记录留存，不直接回写自动化裁决，以避免人工流程破坏原有决策链路的可追溯性。

确认与取消路径是交易状态的最终出口。确认操作只允许在交易状态满足条件时执行；如果交易经过二次质询，则必须先获得通过结果。取消操作允许用户在待处理状态下主动终止交易，取消后系统拒绝后续确认或复核外的状态变更。确认和取消都会写入审计日志，用于还原交易生命周期。

\begin{table}[H]
    \centering
    \caption{各层级在不同流转路径中的职责}
    \label{tab:layer-flow-resp}
    \begin{tabularx}{\linewidth}{|c|Y|Y|Y|}
        \hline
        层级 & 预检路径 & 二次质询路径 & 报告与复核路径 \\
        \hline
        H 层（感知） & 聚合多源数据，封装感知快照 & 提供补充上下文 & 保留原始依据 \\
        \hline
        I 层（认知） & 文本分类，产出风险标签 & 语义分析，产出决策建议 & 生成解释摘要 \\
        \hline
        R 层（规则治理） & 公式计算，阈值路由 & 状态更新，触发报告生成 & 持久化报告，管理复核状态机 \\
        \hline
        D 层（执行） & 执行放行/拦截 & 展示解释面板，执行裁决 & 提供报告入口，展示时间线 \\
        \hline
    \end{tabularx}
\end{table}

\section{感知集成层（H）特征建模}

感知集成层是 HIRD 架构的入口。它不负责判断交易是否欺诈，而是负责把一次转账过程中分散出现的文本、结构化参数和微行为数据整理成后续模块可以使用的快照。对于本文系统来说，H 层的设计重点有两个：一是采集过程不能明显打断用户操作；二是采集结果必须足够规范，能够被 Agent、规则引擎和审计模块共同引用。

\subsection{微行为特征指标体系}

结合第二章对认知负荷和微行为的分析，本文在移动端转账流程中选取五项较容易采集、也较容易解释的行为指标。选择这些指标的原因是，它们都与用户完成转账动作的过程有关，而不是只描述交易最终结果。例如，输入时间变长、停顿增多、反复修改或确认前长时间停留，都可能提示用户正在受到外部信息干扰。当然，这些指标不能单独证明欺诈，只能作为后续综合判断的辅助证据。各指标的判异规则如表\ref{tab:micro-behavior-rules}所示。

\begin{table}[H]
    \centering
    \caption{微行为特征指标及判异规则}
    \label{tab:micro-behavior-rules}
    \begin{tabularx}{\linewidth}{|c|Y|}
        \hline
        指标 & 判异规则 \\
        \hline
        有效输入时长 & 剔除外部干扰后的填写耗时；若较用户基线显著延长则判定异常 \\
        \hline
        非预期停顿次数 & 操作间隔超过 500ms 的停顿数；若次数增多且分布杂乱则判定异常 \\
        \hline
        内容修改频次 & 关键字段的修改次数；频繁修改反映外部指令引导 \\
        \hline
        页面停留时长 & 填完到确认的驻留时间；长时间停留反映决策犹豫 \\
        \hline
        操作节奏方差 & 相邻操作间隔的方差；方差增大反映注意力分散和节奏紊乱 \\
        \hline
    \end{tabularx}
\end{table}

以上五项指标在系统中以特征向量形式使用。本文并不为某一个指标设置绝对拦截条件，而是更关注多项指标的组合偏离。例如，单次页面停留时间变长可能只是用户临时犹豫，但如果同时出现输入时间延长、修改次数增加和操作节奏方差变大，就更值得进入后续语义分析或二次核验。这样的处理方式能够降低单一行为指标带来的误判风险。

\subsection{标准化感知快照生成机制}

移动端采集到的原始事件通常比较碎片化，不能直接交给后端模型或规则引擎使用。H 层需要先对事件进行整理，再形成统一格式的感知快照。本文将这一过程划分为噪声过滤、归一化和快照封装三个步骤。

噪声过滤主要处理设备卡顿、网络延迟、误触和页面渲染造成的异常事件。例如，毫秒级间隔内的连续重复点击更可能来自设备抖动，不应直接计入停顿次数；页面加载造成的等待也需要与用户真实停留区分开。过滤步骤的目标不是追求行为记录越多越好，而是尽量保留与用户决策过程相关的数据。

归一化用于处理不同指标的量纲差异。输入时长通常以秒计，修改频次是离散次数，操作节奏方差则可能落在更小的数值范围内。为避免某一类数值因为量纲较大而在后续计算中占据过高权重，系统采用 Min-Max 方法将连续型指标映射到 $[0,1]$ 区间。

快照封装负责把处理后的行为向量、交易结构化数据和语义文本数据绑定到同一笔交易流水号下，形成标准化感知快照。每笔交易对应一条主要快照，后续 I 层、R 层和审计模块都围绕该快照继续处理。这样可以减少跨层传递时的数据混淆，也方便在风险报告和人工复核中回溯原始依据。

\section{认知决策层（I）详细设计}

认知决策层位于感知快照和规则治理之间，负责把“已经整理好的交易上下文”转化为可供治理层使用的风险线索。它处理的对象包括交易备注、用户补充说明、风险案例匹配结果和部分脱敏后的行为特征。I 层的定位仍然是分析层：它可以给出语义红旗、风险摘要和案例匹配依据，但不直接决定交易放行或拦截。这样的边界设置，能够让 Agent 参与复杂语义判断，同时避免模型输出直接进入资金执行链路。

\subsection{多 Agent 协作架构}

本文将 I 层拆分为四类 Agent，主要是为了降低单个 Agent 同时承担分类、推理、解释和报告生成时带来的不稳定性。初分 Agent 负责快速判断交易属于哪类场景；细分 Agent 在场景范围内匹配更具体的欺诈模式；语义分析 Agent 关注文本中是否存在诱导、威胁、收益承诺等风险表达；报告生成 Agent 则把前面几类结果整理成用户和审核人员能够阅读的结构化说明。各 Agent 的角色定义与输出约束如表\ref{tab:agent-role-matrix}所示。

\begin{table}[H]
    \centering
    \caption{Agent 角色描述与约束矩阵}
    \label{tab:agent-role-matrix}
    \begin{tabularx}{\linewidth}{|c|Y|Y|}
        \hline
        Agent 角色 & 核心任务 & 输出约束 \\
        \hline
        初分 Agent & 对交易场景进行快速归类 & 仅输出场景类别标签，不做深度判定 \\
        \hline
        细分 Agent & 在初分类别下识别具体欺诈模式 & 输出风险评分及匹配案例编号 \\
        \hline
        语义分析 Agent & 解析交易文本中的欺诈意图 & 输出语义红旗标识与推理依据摘要 \\
        \hline
        报告生成 Agent & 将推理过程整理为可读报告 & 仅做说明性输出，不参与业务裁决 \\
        \hline
    \end{tabularx}
\end{table}

四类 Agent 的共同约束是只输出分析结果，不直接调用业务执行工具。前三类 Agent 的输出主要服务于规则治理层，报告生成 Agent 的输出主要服务于解释面板和审计报告。这样设计后，即使某一类 Agent 给出偏高或偏低的判断，也必须经过 R 层的权重、阈值和状态校验后，才能影响最终处置。

\subsection{Agent 协作流程}

四类 Agent 的协作遵循从粗到细的顺序。系统先判断交易大类，再匹配具体风险模式，随后提取语义红旗，最后生成说明性报告。这样安排的好处是，后续 Agent 不需要面对完全开放的文本空间，而是可以在前一阶段已经缩小的场景范围内继续分析。

为更直观地说明 I 层内部的节点协作关系，本文在多 Agent 协作流程图中展示感知快照进入认知层后的处理顺序。该图重点表达四类 Agent 的职责衔接：初分 Agent 先完成场景粗分类，细分 Agent 再结合案例库定位具体风险模式，语义分析 Agent 提取诱导性表达和红旗线索，报告生成 Agent 最后输出结构化解释。图中各节点只产生分析结果，不直接连接交易执行工具，最终处置仍回到 R 层完成。

\begin{figure}
    \centering
    \includegraphics[width=1.0\linewidth]{project_51_3e561d67.png}
    \caption{多 Agent 协作流程图}
    \label{fig:placeholder}
\end{figure}

感知快照从 H 层进入 I 层后，首先由初分 Agent 接收。初分 Agent 依据交易金额、收款信息、设备上下文和文本初步特征，将当前交易归入常规转账、疑似虚假理财、疑似冒充客服、疑似公检法诱导等场景。该阶段只输出类别标签和置信度，不给出最终风险结论。

初分结果产生后，细分 Agent 会在对应类别下检索案例库，并结合外部情报判断当前交易是否接近已知风险模式。例如，若初分结果指向虚假理财，系统会优先匹配收益承诺、限时认购、账户归集等案例特征；若指向冒充客服，则会关注退款验证、订单异常和账户冻结等诱导理由。细分 Agent 输出相似度评分和案例编号，用于保留判断来源。

在细分匹配完成后，语义分析 Agent 进一步处理交易备注、用户补充说明等非结构化文本。该阶段不再只看场景类别，而是提取具体风险表达，例如要求转入“安全账户”、承诺固定收益、强调保密或制造紧迫感等。分析结果以语义红旗列表输出，每个红旗附带简短解释，说明它为什么可能提升交易风险。

前三步完成后，报告生成 Agent 汇总场景类别、案例匹配结果和语义红旗，生成结构化风险说明。报告内容包括风险标题、摘要、风险因子和必要的解释文本。该报告主要用于前端展示、后台复核和审计查询，不参与最终裁决；最终处置仍由 R 层根据评分和阈值完成。

\section{规则治理层（R）决策算法设计}

规则治理层（R 层）是系统实现“风险可控、决策可审计、流程可复核”的关键层。与仅依赖大模型直接给出处置建议的方案不同，本文将大模型、知识库、规则引擎与交易状态机进行统一治理，使语义分析结果只能作为风险证据输入，而不能直接决定资金交易执行。治理层的作用可以概括为三个方面：一是完成多源风险信息的融合计算，二是对交易流程施加状态机约束，三是形成审计记录、风险报告与人工复核闭环。

\subsection{治理层算法设计理念}

治理层算法的设计目标不是单纯提高分类准确率，而是在银行转账场景下实现“可解释、安全、可复核”的风险处置。其设计理念主要包括以下四点。

首先，治理层采用分层风险建模思想，将交易风险拆分为静态风险、行为风险、语义风险和外部情报风险。静态风险主要反映金额、地域、收款人关系、历史转账频率等结构化特征；行为风险反映输入停顿、页面切换、复制粘贴、操作时长等微行为特征；语义风险反映用户备注、补充说明与欺诈话术之间的匹配程度；外部情报风险则来自高风险名单和负面对象筛查。

其次，治理层采用证据收敛而非模型直控策略。系统允许 AI Agent 分析用户文本、提取风险线索、生成解释摘要，但不允许模型直接触发“放行”或“执行转账”。所有交易处置必须经过治理层统一阈值裁决，并满足状态机约束，从而避免模型幻觉、误判或提示注入直接影响资金安全。

再次，治理层强调状态受控。对于预检结果为 block 的交易，确认接口必须拒绝继续执行；对于预检结果为 interrogate 的交易，只有在二次核验结果为 pass\_secondary 时才允许确认；同一笔交易在一次风险核验周期内只允许一次二次提交，防止用户通过重复尝试规避风控。

最后，治理层强调全链可审计。预检开始、预检裁决、二次核验、风险报告生成、人工复核申请与处理等关键节点均写入审计事件表，既支持后续解释，也支持人工复核与责任追踪。

\subsection{治理层风险融合算法}

治理层将识别层输出的风险特征映射为统一风险分，并据此进行三态裁决。设静态风险分为 $S_{static}$，行为风险分为 $G_{behavior}$，语义风险分为 $S_{dev}$，语义匹配强度为 $c_{match}\in [0,1]$。系统首先根据语义命中强度动态调整语义分所占权重，令基础权重为 $\omega_{base}$，则有：

\begin{equation}
\omega_{adj} = \omega_{base} + (1 - \omega_{base})\cdot c_{match}
\end{equation}

在无硬拦截条件时，系统最终风险分定义为：

\begin{equation}
F_{final} = (1 - \omega_{adj})\cdot S_{static} + \omega_{adj}\cdot S_{dev}
\end{equation}

其中，$c_{match}$ 越高，表示文本与高风险场景越接近，此时最终风险分越趋近于语义风险分；反之，若语义证据较弱，则静态风险在总分中占更高权重。

当交易命中高危短语、外部高风险名单且带有阻断提示，或命中治理层定义的强风险规则时，系统进入硬拦截模式，直接令：

\begin{equation}
F_{final} = 1.0
\end{equation}

该机制保证在极端高危条件下，系统不会因加权平均而削弱阻断力度。

在得到最终风险分后，治理层依据阈值进行三态裁决：

\begin{equation}
F_{final}<0.3\Rightarrow \mathrm{pass}
\end{equation}
\begin{equation}
0.3\leq F_{final}\leq 0.8\Rightarrow \mathrm{interrogate}
\end{equation}
\begin{equation}
F_{final}>0.8\Rightarrow \mathrm{block}
\end{equation}

其中，pass 表示允许继续确认，interrogate 表示进入二次核验语义复审路径，block 表示直接拦截交易。

\subsection{治理层算法伪码}

为便于描述治理层的实现过程，本文将其拆分为“预检治理算法”和“二次核验治理算法”两部分。

\begin{verbatim}
算法1：预检治理算法
Input : transfer_request, client_context
Output: decision, risk_level, final_risk, confirmation_token, explain_pack

1 读取账户信息、常用城市、历史收款人关系和近期转账频率
2 调用知识库分类模块，得到 risk_classification
3 调用外部情报筛查模块，得到 external_intelligence
4 根据 risk_classification 计算语义匹配强度 c_match
5 计算静态风险分 S_static
6 计算行为风险分 G_behavior
7 计算语义风险分 S_dev，并判断是否命中硬拦截条件
8 若命中硬拦截，则 final_risk = 1.0
9 否则按动态加权公式融合得到 final_risk
10 根据阈值输出 pass / interrogate / block
11 持久化 pending_transfer、risk_event 和 trace_event
12 生成 explain_pack，返回预检结果
\end{verbatim}

\begin{verbatim}
算法2：二次核验治理算法
Input : confirmation_token, user_reply, context
Output: secondary_decision, final_risk_after_secondary, risk_report

1 根据 confirmation_token 查询待确认交易
2 若交易状态非 pending，则拒绝处理
3 若该交易已提交过二次核验，则拒绝重复提交
4 调用 Agent 对 user_reply 进行语义分析
5 若命中高风险语义红旗，则输出 block_secondary
6 若说明充分且与风险场景不冲突，则输出 pass_secondary
7 若说明不足，则维持 interrogate
8 更新交易的 secondary_decision 与 secondary_risk
9 写入 secondary_decided 审计事件
10 生成并持久化 risk_report
11 写入 risk_report_generated 审计事件
12 返回二次核验结果与解释包
\end{verbatim}

从伪码可以看出，治理层并不是“单次评分后立即完成决策”，而是通过“预检-核验-报告-复核”的多阶段闭环实现风险处置。这也是其区别于单一分类器或单次大模型问答系统的核心特点。

\subsection{治理层时空复杂度分析}

治理层算法由规则融合、知识库匹配、状态机约束和持久化写入共同组成，因此其时间复杂度需要分阶段分析。

在静态风险计算阶段，系统仅涉及金额分段、地域比对、收款人关系判断和频次统计，均为常数级操作，因此该部分时间复杂度为 $O(1)$。在行为风险计算阶段，系统只对有限数量的行为字段进行阈值判断与加权累加，如输入停顿次数、输入时长、页面上下文、复制粘贴次数和应用切换次数等，因此该部分同样为 $O(1)$。

在语义风险计算阶段，系统需要在内置风险场景集中执行文本匹配和向量召回。设风险场景数量为 $m$，向量维度为 $d$，单场景关键词与短语规模为 $k$，则该阶段复杂度可写为：

\begin{equation}
T_{semantic} = O(m \cdot d + m \cdot k)
\end{equation}

当前实现中，风险场景库采用有限内置场景集，$m$ 较小，因此在工程运行中可近似看作常数时间。但从理论上讲，当场景库进一步扩展到更大规模时，该阶段将成为主要计算开销来源，此时需要引入向量索引或分层检索机制以降低检索成本。

在治理状态机阶段，系统只需判断交易是否处于 pending 状态、是否已执行二次核验、是否满足确认条件，这些操作均为常数级，因此复杂度为 $O(1)$。在持久化阶段，系统向数据库写入待确认交易、风险事件、审计事件和风险报告，单次写入记录数固定，因此数据库写入复杂度也可视为 $O(1)$。

综上，治理层在线裁决的总体时间复杂度可表示为：

\begin{equation}
T_{total} = O(1) + O(1) + O(m \cdot d + m \cdot k) + O(1) + O(1)
\end{equation}

即主要受语义风险阶段影响。在当前系统中，由于 $m$ 较小，治理层整体运行时延表现为稳定低延迟。

空间复杂度方面，单次请求主要额外占用当前交易上下文、场景评分结果、解释包结构和审计载荷。若场景集合大小为 $m$，则语义分类阶段需要保留有限候选结果，空间复杂度约为 $O(m)$；其余模块仅维护固定字段和常数级中间变量，因此不改变总体空间复杂度数量级。

\subsection{治理层性能与实验支撑}

从论文结构上看，详细实验过程与测试表格应放在后续测试章节中展开，但在治理层设计部分仍有必要给出简要实验支撑，以证明所提出算法不仅在理论上可行，而且在工程实现中具备可运行性和稳定性。

首先，在自动化回归测试方面，当前系统共收集 66 个后端 PyTest 测试用例，覆盖 API 契约、风险公式、风险场景判定、风险报告生成、人工复核状态机、追踪链路和集成流程等模块，测试结果为 66 项全部通过。这说明治理层的主路径与异常路径均通过了基本正确性验证。

其次，在性能测试方面，系统对预检链路和二次核验链路分别进行了基准测试。结果表明，预检热路径平均响应时间为 43.60ms，P95 响应时间为 78.39ms，说明治理层中的静态风险、行为风险和状态机裁决具有较低额外开销。二次核验热路径平均响应时间为 2360.31ms，P95 为 2767.42ms，该阶段时延明显高于预检，主要原因在于其需要调用语义分析链路并生成风险报告。进一步地，独立的大模型语义路径平均时延为 $5158.52\mathrm{ms}$，P95 为 $5635.89\mathrm{ms}$，这说明二次核验的主要性能瓶颈并不在治理层本身，而在外部语义推理过程。

再次，在一致性与有效性实验方面，系统采用主路径与影子路径对照的方式评估治理输出稳定性。影子模式实验结果显示，决策偏离率为 0，即在测试样本上主路径与影子评估未出现不一致裁决。这说明治理层在当前规则设定下具有较好的决策一致性。

最后，在场景识别实验方面，系统构造并运行了 54 个风险场景样本，其中包含 24 个欺诈样本和 30 个正常样本。实验结果表明，系统在 54 个样本上实现了 50 个精确决策匹配；欺诈召回率达到 1.0000，即所有欺诈样本均被判定为 interrogate 或 block；正常交易直接误拦截率为 0.0000，说明治理层在当前样本集上没有将正常交易直接误判为 block。该结果表明，治理层算法在保持高欺诈召回能力的同时，也较好地控制了对正常交易的过度拦截。

此外，系统还进行了微行为特征消融实验。将输入时长、停顿次数、修改频次、停留时长和操作方差等 H 层微行为字段置为默认值后，当前样本集上的欺诈召回率仍保持为 1.0000，未出现明显下降。该结果说明：在当前实验样本中，语义风险与规则风险特征已足以支撑主要识别能力；而微行为特征的实际价值，更可能体现在高相似边界样本和复杂对抗情形中，这也是后续工作需要进一步扩展样本集验证的方向。

\subsection{本节与后续测试章节的关系说明}

需要说明的是，本节中的性能与实验内容主要承担“算法可行性论证”的作用，因此只保留摘要级结果，用于支撑治理层设计的合理性；而完整的测试过程、测试环境、用例设计、性能表格和消融实验细节，应在后续功能测试与性能测试章节中进一步展开。换言之，本节应当包含性能与实验支撑，但不宜写成完整实验章节，以避免与第四章、第五章形成重复。

\section{业务执行层（D）与全链路审计设计}

业务执行层是 HIRD 架构中最接近用户和交易状态的一层。前三层完成数据采集、语义分析和规则收敛后，D 层只接收 R 层已经确定的指令，并把它转化为页面提示、交易放行、交易拦截、报告展示或复核入口等具体动作。本文将 D 层设计为“只执行、不分析”，目的是避免执行层再次引入模糊判断，也避免 Agent 绕过治理层直接影响交易状态。

D 层的设计有两个重点。第一是状态确定性：系统必须清楚一笔交易当前处于待确认、二次核验、已放行、已拦截、已取消或复核中等哪一种状态。第二是过程可追溯：每一次状态变化都要记录触发来源、时间、输入数据和执行结果。基于这两点，D 层主要承担指令执行、用户交互、MCP 工具调用和审计归档等职责。

\subsection{D 层核心动作}

D 层围绕 R 层输出的三类结果执行动作。对于 pass 结果，D 层保持原转账流程，完成必要的交易记录和审计写入。对于 interrogate 结果，D 层向前端返回二次质询状态，提示用户补充说明或完成核验，再把用户回复交由 I 层和 R 层继续处理。对于 block 结果，D 层阻断交易继续执行，向用户展示风险说明，并保存触发拦截的主要依据。

除自动化处置外，D 层还提供人工复核入口。用户在风险报告页发起复核申请后，系统创建复核记录并进入后台处理队列。审核人员查看交易上下文、风险报告和审计日志后给出处置结论。复核结果作为治理记录保存，不直接覆盖原有自动化裁决，从而保证自动化链路和人工链路都可以被独立追踪。

\subsection{MCP 协议与安全数据交互}

D 层涉及交易执行、风险记录和审计落盘等关键动作，不能通过零散接口随意调用。本文使用 MCP 协议对这些能力进行封装，使上层只能通过明确的工具入口访问底层能力。MCP 在本系统中主要承担三项职责：一是原子工具封装，将风险计算、转账执行、日志写入等能力拆成边界清晰的工具；二是权限管控，每次调用都需要经过身份和范围校验，避免 Agent 接触未授权数据；三是幂等校验，以全局唯一交易流水号作为幂等键，避免重复请求导致重复执行。协议交互的完整时序如图 3.3 所示。

\begin{figure}
    \centering
    \includegraphics[width=1.0\linewidth]{project_51_2ae141e6.png}
    \caption{基于MCP协议的数据交互时序图}
    \label{fig:placeholder}
\end{figure}

\subsection{全链路审计设计}

全链路审计用于回答三个问题：系统看到了什么、为什么做出当前判断、最终执行了什么动作。每笔交易从 H 层感知采集到 D 层指令执行，各处理节点都会生成对应审计事件。整个链路以全局唯一 Transfer ID 贯穿，审计事件链结构如图 3.4 所示。

\begin{figure}
    \centering
    \includegraphics[width=1.0\linewidth]{project_51_0f7171d7.png}
    \caption{全链路追踪流程图}
    \label{fig:placeholder}
\end{figure}

审计日志记录感知快照、Agent 分析结果、R 层评分与阈值路由、D 层执行动作和人工复核状态等信息。后续审核人员可以通过交易流水号回溯各节点输入和输出，判断系统是否存在误判、漏判或状态流转异常。对于本文原型系统而言，审计链路不仅是合规要求，也是验证 HIRD 分层设计是否真正解耦的重要依据。

\section{本章小结}

本章围绕手机银行转账场景，对系统需求和总体架构进行了设计。首先，从功能性和非功能性两个角度梳理了系统需要解决的问题，包括多模态感知、语义识别、权限约束、分级治理和审计追踪。其次，补充 Use-case 建模并提出 HIRD 四层架构，说明感知、认知、治理和执行之间的数据流转关系。随后，本章设计了 H 层的五项微行为指标和感知快照机制，说明 I 层中四类 Agent 的协作方式，并给出 R 层的动态权重、分级路由、算法伪码、复杂度和实验支撑。最后，本章说明 D 层如何完成业务执行、MCP 安全交互和全链路审计。

通过上述设计，本文把 AI Agent 的语义分析能力限制在认知层，把最终交易处置交给规则治理层和业务执行层完成。这样的分层方式既保留了大模型处理复杂文本的优势，也使交易执行、人工复核和审计记录保持确定性。第四章将在本章设计基础上说明系统的具体工程实现。
\chapter{系统实现与功能验证}

本章在第三章系统设计的基础上，对手机银行语义欺诈风控原型系统的工程实现与功能验证过程进行说明。系统实现围绕移动端微行为采集、后端风险预检、AI Agent 语义分析、MCP 工具调用、风险报告生成和人工复核流程展开。与第三章偏重架构设计不同，本章重点说明系统在代码层面的模块划分、接口流转方式和功能测试结果。

\section{系统开发环境与技术栈选型}

本文系统采用前后端分离的实现方式。移动端负责转账界面、风险提示和微行为采集；后端负责交易预检、二次核验、风险报告、人工复核、MCP 工具调用和审计记录。系统开发与测试环境如表\ref{tab:dev-env}所示。

\begin{table}[H]
    \centering
    \caption{系统开发与运行环境配置表}
    \label{tab:dev-env}
    \begin{tabularx}{0.75\linewidth}{|c|Y|}
        \hline
        配置项 & 配置内容 \\
        \hline
        操作系统 & Windows 11 \\
        \hline
        后端语言 & Python 3.11 \\
        \hline
        后端框架 & FastAPI \\
        \hline
        测试框架 & PyTest \\
        \hline
        移动端框架 & Flutter \\
        \hline
        语义分析模块 & AI Agent、RAG 检索、大语言模型接口 \\
        \hline
        数据交互方式 & MCP 工具封装、Trace ID 审计追踪 \\
        \hline
    \end{tabularx}
\end{table}

围绕 HIRD 架构的实现需求，系统选取的核心技术栈如表\ref{tab:tech-stack}所示。各技术组件并不是简单堆叠，而是分别对应系统中的具体职责：Flutter 负责移动端交互与行为采集，FastAPI 负责后端异步接口编排，MCP 负责工具调用边界和审计链路，AI Agent 与 RAG 负责语义风险分析，PyTest 则用于功能回归测试和结果复现。

\begin{table}[H]
    \centering
    \caption{系统核心技术栈选型表}
    \label{tab:tech-stack}
    \small
    \begin{tabularx}{\linewidth}{|Y|Y|Y|}
        \hline
        架构层级 / 模块 & 技术组件 & 核心作用 \\
        \hline
        感知集成层 & Flutter & 支撑移动端转账界面、事件监听和微行为特征采集 \\
        \hline
        业务网关 & FastAPI & 提供异步接口、参数校验和交易状态流转能力 \\
        \hline
        认知决策层 & AI Agent + RAG & 对交易备注、用户补充说明和风险案例进行语义分析 \\
        \hline
        规则治理层 & 风险评分与阈值路由模块 & 将语义风险线索和规则分值收敛为确定性处置结果 \\
        \hline
        业务执行层 & MCP 工具链 & 封装风险计算、审计写入、交易执行和幂等校验 \\
        \hline
        测试验证层 & PyTest & 对接口、状态机、风险报告、MCP 追踪和场景样本进行自动化验证 \\
        \hline
    \end{tabularx}
\end{table}

\section{核心模块实现}

\subsection{基于 MCP 协议的安全数据中枢实现}

MCP 数据中枢的实现目标，是让 AI Agent 在分析交易风险时只能通过受控工具访问必要信息，而不是直接接触底层业务接口。系统将风险计算、交易状态变更、审计记录写入等能力封装为独立工具，并在工具内部加入权限校验、幂等检查和日志记录。这样做可以降低模型输出直接影响交易状态的风险，也便于后续追踪每一次工具调用的来源和结果。

在系统实现中，MCP 工具以交易流水号和 Trace ID 作为关键关联字段。交易进入风险预检流程后，后端会为其生成唯一标识，后续 Agent 分析、规则治理、二次核验、风险报告和人工复核均围绕该标识展开。对于重复调用或并发重放请求，MCP 工具层通过幂等校验避免重复执行，从而保证交易状态不会因为接口重试而出现二次漂移。

\subsection{移动端微行为感知模块实现}

移动端感知模块对应 HIRD 架构中的 H 层，主要负责在转账流程中采集用户操作过程数据。本文关注的是输入时长、非预期停顿、修改次数、页面停留和操作节奏方差等指标，而不是保存完整的用户输入过程。这样可以在获取风险判断线索的同时，减少对用户隐私和界面体验的影响。

实现上，系统利用 Flutter 的事件回调机制记录转账界面中的关键交互事件，包括字段输入、内容修改、页面停留和提交动作。用户提交交易时，移动端将行为统计结果与交易金额、收款信息、备注文本等字段一起封装为风险感知快照。该快照作为后端预检和语义分析的输入，为后续规则治理提供统一数据基础。

\subsection{后端认知决策与异步编排实现}

后端实现围绕预检热路径和语义冷路径两条流程展开。预检热路径主要处理本地规则、结构化字段和状态记录，要求尽快返回；语义冷路径涉及 AI Agent、RAG 检索和大语言模型接口，耗时更长，需要通过异步挂起和状态机配合。

感知快照到达后端后，系统先执行风险预检，得到初步处置结果。若交易风险较低，则返回放行结果；若命中高风险规则，则进入拦截流程；若处于中间风险区间，则触发二次核验。用户提交补充说明后，后端结合本地规则与 Agent 语义分析结果再次判断，最终由规则治理层收敛为放行、二次核验或拦截等确定性结果。

\section{系统功能测试与验收验证}

为验证系统主要功能是否能够闭环，本文使用 PyTest 自动化测试框架对后端接口、风险公式、二次核验、人工复核、风险报告和 MCP 审计链路进行验证。测试数据以场景化构造样本为主，包括常规转账、虚假理财、冒充公检法、冒充客服和异常停顿等类型。

本次功能测试的真实执行结果来自项目当前代码版本。通过执行 \texttt{pytest --collect-only -q}，系统共采集到 66 个后端 PyTest 测试用例；已有全量执行日志记录结果为 66 passed in 53.25s，说明当前版本中的后端功能链路均通过自动化回归验证。需要说明的是，该 66 个用例只统计 \texttt{server/tests} 目录下的后端自动化测试，不包含 Flutter 前端 \texttt{widget\_test.dart} 中的 23 个 widget 测试，避免前后端测试口径混合。

\begin{table}[H]
    \centering
    \caption{系统功能测试用例与结果汇总}
    \label{tab:func-test-cases}
    \small
    \begin{tabularx}{\linewidth}{|c|Y|Y|Y|c|}
        \hline
        编号 & 测试目标 & 覆盖脚本/接口 & 关键断言 & 结果 \\
        \hline
        FT-01 & 转账预检主流程 & \texttt{test\_api.py} 中的 flow & 返回决策、风险分、解释包结构完整 & 通过 \\
        \hline
        FT-02 & 语义分类高危短语识别 & \texttt{test\_api.py} 中高危短语识别相关用例 & 命中高危短语并提升风险级别 & 通过 \\
        \hline
        FT-03 & 否定语义鲁棒性 & \texttt{test\_api.py} 中否定语义相关用例 & 否定表达不应误触发高危命中 & 通过 \\
        \hline
        FT-04 & 外部情报联动 & \texttt{test\_api.py} 中的 report & 命中外部名单时状态与解释节点一致 & 通过 \\
        \hline
        FT-05 & Block 后禁止确认 & \texttt{test\_api.py} 中的 confirm & 被拦截交易确认请求返回 400 & 通过 \\
        \hline
        FT-06 & 二次核验放行路径 & \texttt{test\_api.py} 中 \texttt{pass\_secondary} 相关用例 & 返回 \texttt{pass\_secondary} 且解释包完整 & 通过 \\
        \hline
        FT-07 & 二次核验拦截路径 & \texttt{test\_api.py} 中 \texttt{block\_secondary} 相关用例 & 返回 \texttt{block\_secondary} 且红旗信息非空 & 通过 \\
        \hline
        FT-08 & 人工复核状态机 & \texttt{test\_manual\_review.py} 系列 & \texttt{submitted}、\texttt{in\_review}、\texttt{closed} 合法流转 & 通过 \\
        \hline
        FT-09 & 风险报告生成与查询 & \texttt{test\_risk\_report.py} 系列 & secondary-check 后可持久化并查询报告 & 通过 \\
        \hline
        FT-10 & Trace 与 MCP 一致性 & \texttt{test\_trace\_mcp.py} 系列 & 审计事件可追溯，查询结果与仓储一致 & 通过 \\
        \hline
    \end{tabularx}
\end{table}

为进一步说明测试覆盖范围，本文按测试脚本统计自动化用例数量，结果如表\ref{tab:auto-test-stat}所示。

\begin{table}[H]
    \centering
    \caption{自动化功能测试统计}
    \label{tab:auto-test-stat}
    \begin{tabularx}{\linewidth}{|Y|c|c|}
        \hline
        测试文件 & 用例数 & 通过数 \\
        \hline
        \path{server/tests/test_api.py} & 24 & 24 \\
        \hline
        \path{server/tests/test_formula.py} & 6 & 6 \\
        \hline
        \path{server/tests/test_manual_review.py} & 7 & 7 \\
        \hline
        \path{server/tests/test_patch2_secondary_followup.py} & 2 & 2 \\
        \hline
        \path{server/tests/test_risk_report.py} & 4 & 4 \\
        \hline
        \path{server/tests/test_risk_scenarios.py} & 14 & 14 \\
        \hline
        \path{server/tests/test_trace_mcp.py} & 6 & 6 \\
        \hline
        \path{server/tests/test_v3_integration_flows.py} & 3 & 3 \\
        \hline
        总计 & 66 & 66 \\
        \hline
    \end{tabularx}
\end{table}

从测试分布可以看出，后端自动化测试并不只覆盖单一接口，而是覆盖了转账预检、风险公式、二次核验、人工复核、风险报告、MCP 审计追踪和集成流程等多个模块。其中，\texttt{test\_api.py} 用例数量最多，主要验证核心 API 流程；\texttt{test\_risk\_scenarios.py} 负责低风险放行、中风险质询、高风险拦截、否定语义和典型诈骗场景验证；\texttt{test\_trace\_mcp.py} 则用于确认审计链路与 MCP 工具调用的一致性。上述结果说明，当前原型系统在既定后端测试用例范围内能够完成主要业务闭环。

\section{移动端界面展示与核心交互流程}

系统移动端界面用于承载普通转账、风险提示、二次质询、账单查询、AI 助手和风险报告等功能。前端不重新计算风险，也不直接决定交易处置，而是根据后端返回的状态、解释信息和报告数据完成页面展示。这样可以保证风控逻辑集中在后端治理流程中，移动端主要承担交互入口和结果呈现职责。

首页界面是用户进入原型系统后的主要入口，集中展示账户概览、常用功能和转账入口。该页面对应后端的基础账户与功能导航能力，不触发复杂风控计算，但它为后续转账链路提供统一入口。用户从首页进入转账功能后，系统才开始采集本次交易过程中的输入时长、页面停留和操作节奏等微行为数据。

\begin{figure}[H]
    \centering
    \includegraphics[width=0.75\linewidth]{project_51_803335f7.png}
    \caption{首页界面}
    \label{fig:placeholder}
\end{figure}

转账功能界面用于填写收款方、转账金额、交易备注等关键字段，也是 H 层微行为采集的主要发生位置。用户输入和修改字段时，前端记录停顿次数、修改频次、页面停留时长等指标，并在提交交易时与交易参数一起封装为感知快照。后端收到请求后进入风险预检流程，返回 pass、interrogate 或 block 三类状态。


\begin{figure}[H]
    \centering
    \includegraphics[width=0.75\linewidth]{project_51_28d96d47.png}
    \caption{转账功能界面}
    \label{fig:placeholder}
\end{figure}

二次质询界面用于处理预检结果为 interrogate 的交易。当前端收到该状态后，会提示用户补充说明转账背景或资金用途，并将用户输入提交到二次核验接口。该页面对应后端的 secondary-check 流程，用户补充文本会作为新的语义分析材料进入 Agent 分析和规则治理流程；若回复能够解释交易合理性，系统可继续放行，否则进入拦截和报告生成流程。

\begin{figure}[H]
    \centering
    \includegraphics[width=0.75\linewidth]{project_51_ce91b243.png}
    \caption{二次质询界面}
    \label{fig:placeholder}
\end{figure}

账单明细界面用于展示交易记录、交易状态和风险处置结果。对于已放行、已拦截或进入复核流程的交易，用户可以通过该页面查看对应流水和状态变化。该页面与后端交易状态机和审计记录查询能力对应，能够把预检、二次核验、确认、取消和复核等节点以用户可理解的方式呈现出来。

\begin{figure}[H]
    \centering
    \includegraphics[width=0.75\linewidth]{project_51_3654f6f6.png}
    \caption{账单明细界面}
    \label{fig:placeholder}
\end{figure}

AI 助手界面用于展示系统对风险原因的解释性提示，重点呈现语义红旗、风险摘要和用户可理解的处置建议。该页面不承担业务裁决职责，只读取后端生成的解释包或风险报告摘要。通过将 Agent 分析结果转化为较清晰的前端说明，系统能够降低用户对拦截或二次核验结果的理解成本。


\begin{figure}[H]
    \centering
    \includegraphics[width=0.75\linewidth]{project_51_fb6ce5d7.png}
    \caption{AI 助手界面}
    \label{fig:placeholder}
\end{figure}

风险报告界面用于展示被拦截或进入风险处置交易的完整说明。报告内容包括风险标题、综合风险等级、风险摘要、风险因子列表和人工复核入口。该页面对应后端风险报告生成与查询接口，也关联 Trace ID 审计链路。若用户对处置结果存在异议，可以从该页面提交人工复核申请，审核人员后续依据交易上下文、风险报告和审计事件给出复核结论。

\begin{figure}[H]
    \centering
    \includegraphics[width=0.75\linewidth]{project_51_20859ddf.png}
    \caption{风险报告界面}
    \label{fig:placeholder}
\end{figure}

\section{本章小结}

本章围绕系统工程实现和功能验证进行了说明。系统采用 Flutter 与 FastAPI 完成前后端分离实现，通过 AI Agent 和 RAG 支撑语义分析，通过 MCP 工具链约束风险计算、交易执行和审计写入边界。移动端界面部分展示了首页、转账、二次质询、账单明细、AI 助手和风险报告等功能页面，说明前端如何承接后端风控状态并完成用户交互。在功能验证方面，当前项目共采集到 66 个后端 PyTest 自动化测试用例，已有执行日志结果为 66 passed in 53.25s。测试覆盖接口流程、风险公式、二次核验、人工复核、风险报告、MCP 追踪和集成流程等模块。结果表明，在当前测试范围内，系统能够完成手机银行语义欺诈风控原型的主要功能闭环。

\chapter{实验测试与分析}

本章在第四章功能验证的基础上，进一步对系统的场景识别效果、关键链路性能、影子模式一致性和消融实验结果进行分析。与第四章关注“功能是否跑通”不同，本章重点关注“系统在既定场景样本和测试脚本下表现如何”。需要说明的是，本章实验基于当前项目代码与场景化测试数据展开，实验结论适用于本文构造的原型系统测试条件，不等同于真实银行生产环境中的最终效果。

\section{实验环境与执行方法}

实验在本地开发环境中执行，后端运行环境为 Python 3.11，测试框架为 PyTest。实验过程中主要使用如下命令完成测试采集、回归验证、性能基准测试和场景指标统计：

\begin{verbatim}
server\.venv\Scripts\python -m pytest --collect-only -q
server\.venv\Scripts\python -m pytest server\tests -q
server\.venv\Scripts\python server\tests\benchmark_v3_baseline.py
server\.venv\Scripts\python server\tests\benchmark_v3_llm_path.py
server\.venv\Scripts\python server\tests\benchmark_v3_precheck_shadow.py
server\.venv\Scripts\python server\tests\analyze_risk_scenarios_metrics.py
\end{verbatim}

其中，\texttt{pytest --collect-only -q} 用于统计后端测试用例数量，结果显示共采集 66 个 PyTest 用例；\texttt{pytest server/tests -q} 用于执行后端回归测试，已有日志记录结果为 66 passed in 53.25s。\texttt{benchmark\_v3\_baseline.py} 用于统计 \texttt{precheck} 热路径和 \texttt{secondary-check} 路径耗时；\texttt{benchmark\_v3\_llm\_path.py} 用于统计真实 LLM 语义链路耗时；\texttt{benchmark\_v3\_precheck\_shadow.py} 用于验证影子模式一致性；\texttt{analyze\_risk\_scenarios\_metrics.py} 用于计算 54 个场景样本的决策匹配度、欺诈召回率、正常交易直接误拦截率和消融结果。

\section{测试用例来源与覆盖范围}

为回应“测试用例是什么、如何测试、验证什么”的问题，本文将测试用例分为功能回归用例、场景识别用例和性能观测用例三类。功能回归用例来自 \texttt{server/tests} 目录下的 PyTest 脚本，主要用于验证接口契约、状态流转和数据持久化是否符合设计；场景识别用例来自 54 个带有人工标注的风险样本，其中包含 24 个欺诈样本和 30 个正常样本，主要用于观察系统对正常交易、明确欺诈和边界模糊交易的判定效果；性能观测用例来自基准测试脚本，主要用于统计不同链路的响应时延和一致性指标。

功能回归测试覆盖系统主要业务闭环。其测试对象包括转账预检主流程、风险公式计算、二次核验通过与拦截、风险报告生成与查询、人工复核状态机、Trace 与 MCP 审计一致性，以及跨模块集成流程。测试执行时，PyTest 会构造不同交易请求、用户补充说明和状态变更条件，调用后端接口并断言返回状态、风险等级、审计事件和持久化结果是否符合预期。因此，66 个后端自动化用例并不是简单的接口连通性测试，而是覆盖了从预检到复核的关键状态路径。前端 \texttt{widget\_test.dart} 中还包含首页、弹窗、风险报告中心、人工复核界面、等待提示和前端数据解析等 widget 测试，但本文第五章中的该统计值仅指后端 PyTest 结果。

场景识别测试用于说明系统在不同风险样本上的判定来源。本文使用的 54 个样本包含普通转账、虚假理财、冒充客服、冒充公检法、异常停顿和边界模糊场景。测试脚本逐条输入样本，记录系统输出的 pass、interrogate 或 block 结果，并与人工标注进行比较，最终统计 exact match、欺诈召回率和正常交易直接误拦截率。该类测试与第四章的 UI 展示不同：UI 截图用于说明功能页面已经实现，场景识别用例用于验证后端风控逻辑在样本集上的输出是否符合预期。

性能观测测试覆盖 precheck 热路径、secondary-check 路径、LLM 语义路径和影子模式一致性。precheck 热路径用例模拟普通交易连续预检，验证系统是否能够在 100ms 设计目标内返回；secondary-check 用例模拟需要用户补充说明的交易，验证二次核验链路是否满足异步交互要求；LLM 语义路径用例单独观察外部模型分析耗时，用于说明其不适合作为所有交易的同步默认路径；影子模式用例则在同一批输入上同时执行主路径与 shadow 路径，验证两者裁决是否存在偏离。论文后续测试结果均引用这些已有脚本和日志输出，不额外编造未执行的实验数据。

\section{场景识别与决策一致性实验}

场景识别实验用于验证系统在不同风险样本上的决策效果。测试样本共 54 个，其中欺诈样本 24 个、正常样本 30 个，覆盖正常转账、语义欺诈、边界模糊案例和需要二次核验的交易场景。系统输出结果与人工标注结果进行比较，统计 exact match、欺诈召回率和正常交易直接误拦截率。实验结果如表\ref{tab:scenario-consistency}所示。

\begin{table}[H]
    \centering
    \caption{场景识别与一致性实验结果}
    \label{tab:scenario-consistency}
    \begin{tabularx}{\linewidth}{|c|Y|c|}
        \hline
        实验项 & 指标定义 & 结果 \\
        \hline
        54 样本决策匹配度 & exact decision match / total & 50/54 \\
        \hline
        欺诈召回率 & fraud 样本中判为 interrogate 或 block 的比例 & 1.0000 \\
        \hline
        正常交易直接误拦截率 & normal 样本中被判为 block 的比例 & 0.0000 \\
        \hline
        影子模式偏离率 & 主路径与 shadow 路径决策不一致比例 & 0.0000 \\
        \hline
    \end{tabularx}
\end{table}

从结果看，系统在 54 个样本中有 50 个样本的决策结果与人工标注完全一致，决策匹配度为 50/54。欺诈召回率为 1.0000，说明测试集中的欺诈样本均被系统识别为需要二次核验或拦截；正常交易直接误拦截率为 0.0000，说明正常样本没有被系统直接判为 block。影子模式偏离率为 0.0000，说明在当前测试条件下，主路径与 shadow 路径输出保持一致，预检决策具有较好的稳定性。

需要注意的是，50/54 的 exact match 并不意味着系统存在严重漏判。该指标要求系统输出与人工标注完全一致，而欺诈召回率和正常交易直接误拦截率分别关注安全侧和用户体验侧。从实验结果看，当前系统在安全侧保持了较强的召回能力，同时没有直接误拦截正常样本。但由于样本规模仍然有限，该结果仍需在更丰富的真实或半真实场景数据上继续验证。

\section{关键链路性能实验}

性能测试重点关注三条链路：预检热路径、二次核验路径和 LLM 语义路径。预检热路径用于处理常规交易，要求尽量快速返回；二次核验路径用于处理中风险交易，需要接收用户补充说明并重新判断；LLM 语义路径涉及外部模型调用，耗时受网络、模型响应速度和输入长度影响较大。测试结果如表\ref{tab:key-path-performance}所示。

\begin{table}[H]
    \centering
    \caption{关键链路性能测试结果}
    \label{tab:key-path-performance}
    \begin{tabularx}{0.75\linewidth}{|Y|c|c|c|}
        \hline
        测试项 & 平均时延 (ms) & P95 时延 (ms) & 轮次 \\
        \hline
        precheck 冷启动 & 24.09 & 24.09 & 1 \\
        \hline
        precheck 热路径 & 43.60 & 78.39 & 50 \\
        \hline
        secondary-check 冷启动 & 3192.95 & 3192.95 & 1 \\
        \hline
        secondary-check 热路径 & 2360.31 & 2767.42 & 30 \\
        \hline
        LLM 语义路径 & 5158.52 & 5635.89 & 5 \\
        \hline
    \end{tabularx}
\end{table}

预检热路径平均时延为 43.60 ms，P95 时延为 78.39 ms，能够满足系统在第三章中提出的预检链路 100 ms 内返回的设计目标。该路径主要执行本地规则计算、状态检查和必要的解释信息封装，不主动调用大语言模型，因此时延相对稳定。

二次核验热路径平均时延为 2360.31 ms，P95 时延为 2767.42 ms。与预检路径相比，该路径耗时明显增加，原因在于系统需要处理用户补充说明，并结合语义分析和状态机逻辑重新生成处置结果，适合作为中风险交易的异步确认流程。

LLM 语义路径平均时延为 5158.52 ms，P95 时延为 5635.89 ms，明显高于预检和二次核验路径。这说明大模型调用不适合作为所有交易的默认同步步骤。本文系统采用冷热分流策略：普通交易优先走预检热路径，只有在本地规则和二次核验无法充分判断时，才进入 LLM 语义路径。该策略能够在保持常规交易体验的同时，为复杂语义欺诈场景提供更深层的分析能力。

\section{性能与精度测试用例设计}

为了使第五章实验结果具备可复现性，本文将性能与精度测试拆分为若干测试用例。各用例的输入样本、测试步骤和评价指标如表\ref{tab:perf-accuracy-cases}所示。

\begin{table}[H]
    \centering
    \caption{性能与精度测试用例表}
    \label{tab:perf-accuracy-cases}
    \small
    \begin{tabularx}{\linewidth}{|c|Y|Y|Y|Y|}
        \hline
        编号 & 测试目标 & 输入样本 & 测试步骤 & 评价指标 \\
        \hline
        PT-01 & 预检热路径时延 & 普通转账请求样本 & 连续调用预检接口并记录每次耗时 & 平均响应时间、P95 响应时间 \\
        \hline
        PT-02 & 二次核验路径时延 & 需要二次核验的交易样本 & 提交用户补充说明并记录端到端耗时 & 平均响应时间、P95 响应时间 \\
        \hline
        PT-03 & 大模型调用路径时延 & 需要语义分析的风险文本样本 & 调用 Agent 或 LLM 分析链路并记录耗时 & 平均响应时间、P95 响应时间 \\
        \hline
        PT-04 & 场景识别精度 & 54 个标注风险场景样本 & 对每个样本执行系统判定并与标注比较 & exact match、欺诈召回率 \\
        \hline
        PT-05 & 正常交易直接误拦截 & 标注为 normal 的样本集 & 执行系统判定并统计 block 数量 & 正常交易直接误拦截率 \\
        \hline
        PT-06 & 影子模式一致性 & 同一批预检样本 & 主路径和 shadow 路径分别执行并比较输出 & 决策偏离率 \\
        \hline
    \end{tabularx}
\end{table}

上述测试用例覆盖了系统性能和识别效果的主要观察维度。其中，PT-01 至 PT-03 关注不同链路的响应时间，PT-04 和 PT-05 关注场景识别质量，PT-06 关注同一输入下主路径与影子路径是否保持一致。通过将实验结果与测试用例对应，可以避免只给出指标而缺少测试来源的问题。

\section{消融实验与分析}

为观察 H 层微行为特征对当前测试集的影响，本文进行了 H 层特征消融实验。实验做法是在保持其他输入条件不变的情况下，去除输入时长、停顿频率、修改次数、页面停留和操作节奏方差等 5 项微行为特征，再对同一批场景样本重新计算欺诈召回率。实验结果如表\ref{tab:h-layer-ablation}所示。

\begin{table}[H]
    \centering
    \caption{H 层微行为特征消融实验结果}
    \label{tab:h-layer-ablation}
    \begin{tabularx}{0.75\linewidth}{|Y|c|c|}
        \hline
        配置 & 欺诈召回率 & 召回率变化 \\
        \hline
        完整特征（baseline） & 1.0000 & $-$ \\
        \hline
        去除 H 层 5 项微行为特征后 & 1.0000 & 0.0000 \\
        \hline
    \end{tabularx}
\end{table}

从结果看，在当前 54 个场景样本上，去除 H 层 5 项微行为特征后，欺诈召回率仍为 1.0000，召回率没有发生变化。这一结果说明，当前测试集中的欺诈样本多数已经具有较明显的文本或结构化风险特征，即使不依赖微行为特征，系统仍然能够识别这些样本。

该实验结果并不意味着微行为特征没有价值，而是说明当前测试样本还不足以充分体现 H 层特征的独立贡献。微行为特征更可能在“交易备注较正常，但用户操作节奏异常”的隐蔽场景中发挥作用。因此，后续测试需要增加弱文本风险样本、行为异常样本和规则与语义结果冲突的边界样本，以进一步验证 H 层特征对系统识别能力的增益。

\section{本章小结}

本章基于当前项目代码和真实测试输出，对系统功能、场景识别、关键链路性能、影子模式一致性和消融实验进行了分析。功能回归测试中，系统共采集 66 个后端 PyTest 用例，已有全量执行日志结果为 66 passed in 53.25s。场景识别实验中，54 个样本包含 24 个欺诈样本和 30 个正常样本，exact match 为 50/54，欺诈召回率为 1.0000，正常交易直接误拦截率为 0.0000。影子模式测试中，主路径与 shadow 路径的决策偏离率为 0.0000。

性能方面，precheck 热路径平均时延为 43.60 ms，P95 时延为 78.39 ms；secondary-check 热路径平均时延为 2360.31 ms，P95 时延为 2767.42 ms；LLM 语义路径平均时延为 5158.52 ms，P95 时延为 5635.89 ms。实验结果说明，系统采用冷热分流是必要的：常规交易可通过本地预检快速返回，复杂语义风险则通过二次核验和 LLM 路径进一步分析。消融实验显示，当前测试集尚未充分体现 H 层微行为特征的独立贡献，这也是后续样本扩充和实验设计需要重点改进的方向。

\chapter{结论与展望}

本文围绕手机银行转账场景中的语义欺诈防控问题，完成了从需求分析、架构设计、核心模块实现到自动化测试验证的完整研究过程。与传统依赖金额、地域、设备状态等结构化字段的风控方式相比，本文更关注用户在被诱导或胁迫状态下产生的操作节奏变化，并将微行为特征、交易文本语义和确定性规则治理结合起来，形成一套基于 AI Agent 与 MCP 协议的 HIRD 分层风控原型系统。需要说明的是，本文实验主要依托模拟场景数据和项目自动化测试展开，相关结论主要成立于本文构造的测试条件之内。

\section{主要研究结论}

结合系统实现和第五章测试结果，本文得到以下结论：

1、HIRD 分层架构能够支撑手机银行转账场景下的风险闭环。本文将系统划分为感知集成层、认知决策层、规则治理层和业务执行层，分别承担数据采集、语义分析、决策收敛和动作执行职责。在工程实现中，移动端负责生成微行为快照，后端负责预检、二次核验、风险报告和人工复核等流程编排。第五章测试结果显示，系统全部 66 个后端 PyTest 用例均通过自动化验证，通过率为 100\%，说明各模块在当前实现范围内能够完成预期业务闭环。

2、双路径决策机制在实时性和深度分析之间取得了可验证的平衡。对于常规交易，系统优先执行预检热路径，由本地规则和结构化字段完成快速判定；对于需要进一步确认的中高风险交易，系统再进入二次核验或 Agent 语义分析路径。性能测试表明，预检热路径平均响应时间为 43.60ms，P95 为 78.39ms，满足 100ms 内完成预检的设计目标；二次核验路径平均响应时间为 2360.31ms，P95 为 2767.42ms，满足异步交互场景下 3s 内响应的约束。大模型调用路径平均耗时为 5158.52ms，虽然明显高于本地规则路径，但通过异步挂起和状态机流转，系统仍能保持业务流程的连续性。

3、场景化测试验证了系统对典型欺诈样本的识别能力。本文基于 54 个核心场景样本，对常规转账、虚假理财、冒充公检法、冒充客服和异常停顿等场景进行测试。实验结果显示，系统精确匹配率为 92.6\%，欺诈召回率为 100\%，正常交易直接误拦截率为 0\%。这些结果说明，在当前构造的测试集中，系统能够较稳定地区分正常交易与典型语义欺诈交易。同时，影子模式测试中预检路径判定偏离率为 0.0\%，表明在相同输入条件下，系统主路径与镜像路径能够保持一致输出。

4、MCP 协议为工具调用和审计记录提供了统一边界。系统将风险计算、转账执行、审计记录等能力封装为 MCP 原子工具，并通过 Trace ID 串联预检、Agent 分析、规则治理和执行结果。该设计使 Agent 主要承担语义线索分析，不直接操作资金执行动作；最终业务动作仍由确定性代码逻辑接管。通过这种方式，系统在利用大模型语义理解能力的同时，为后续审计、复核和问题定位保留了较清晰的调用链路。

\section{本文创新点总结}

本文工作主要体现在以下三个方面：

1、将微行为快照引入手机银行转账风控链路。本文没有只依赖金额、设备、IP 等静态字段，而是从用户操作过程入手，将输入时长、停顿频率、修改次数、页面停留和操作节奏方差等特征封装为风险快照。该设计为识别“交易参数正常但用户可能被诱导”的场景提供了额外观察维度。

2、设计了分析权与执行权分离的 HIRD 架构。本文将 AI Agent 放置在认知决策层，主要负责文本语义解析、风险线索提取和报告生成；规则治理层负责将语义结果与硬规则评分进行整合；业务执行层只处理放行、二次核验、拦截和人工复核等确定性动作。这种结构降低了 Agent 直接影响资金流转的风险，也使系统边界更容易被测试和审计。

3、在原型系统中落地了 MCP 风控工具链。系统围绕 MCP 协议封装风险计算、转账执行和审计记录等工具，并结合幂等校验、状态机流转和 Trace ID 追踪机制，形成了一条可复查的调用链路。相比普通接口调用，该方式更适合需要权限边界、日志追踪和复核依据的智能风控场景。

\section{研究不足与未来展望}

本文虽然完成了原型系统开发和自动化测试验证，但仍存在以下不足。

首先，当前实验主要依托 54 个核心场景样本和项目自动化测试展开，尚未接入真实银行脱敏流水。场景化样本能够覆盖典型诈骗话术和异常操作节奏，但难以完整反映真实业务中的长尾用户习惯、复杂交易关系和突发并发压力。因此，本文得到的精确匹配率、召回率、正常交易直接误拦截率和时延数据，仍需要在更大规模、更贴近真实业务的数据集上继续验证。

其次，第五章消融实验显示，在当前 54 个样本中，移除 H 层微行为特征或将 $\omega_{adj}$ 设置为 0 后，系统精度指标未出现明显下降。该现象说明当前样本中的欺诈文本特征较为显性，结构化规则和语义分析已经能够覆盖大部分风险判断。换言之，本文尚未充分构造出能够体现微行为特征和动态权重机制独立贡献的隐蔽样本，这也是后续实验设计需要重点补强的部分。

再次，系统目前主要围绕手机银行转账业务展开，尚未覆盖理财申购、贷款审批、账户变更等更多金融操作。不同业务的风险目标、用户行为路径和合规要求并不相同，现有 HIRD 架构虽然具备扩展空间，但 MCP 工具定义、阈值设置和前端交互方式仍需要结合具体业务重新设计。

后续研究可从三个方向继续推进：第一，扩展真实或半真实的风险案例库，增加弱文本特征、长尾诈骗话术和复杂微行为模式样本，使多模态感知与动态权重机制的贡献能够被更充分地验证；第二，结合更多用户历史行为与设备环境特征，优化 Agent 语义分析与规则治理之间的协同方式，降低对显性关键词的依赖；第三，在更接近生产环境的并发条件下继续测试 MCP 工具链的权限控制、幂等防护、异常恢复和审计追踪能力，使系统逐步从课程原型走向更接近真实金融业务的工程形态。

\end{document}
